\section{问题二模型建立与求解}
\label{sec:question-two}

问题一把每个子图看作独立Task，子图边界可能带来写回、读入和激活等待。问题二仍提交操作归属、分核和核内子图序列，但同一核心上的子图合成一个核级Task。数据可在该核心的L1或UB内继续使用，跨核心的数据仍经DDR传送。因此，本问的决策不只是减少跨核边：把消费者集中在同一核心可以节省读入，也可能让几个大张量同时留在片上，随后发生换出与恢复。本节依次建立核心级结构搬运、物理副本容量、有限候选求解和逐图结果之间的关系。

\subsection{数据预处理与核心驻留域}
\subsubsection{核心驻留域划分}
给定第\ref{sec:question-one}节的计划$X=(g,a,\sigma)$，先由\texttt{node\_to\_subgraph}取得操作所属子图$g(v)$，再从\texttt{core\_schedules}取得子图所在核心$a(g(v))$及同核先后次序。原计算图中的COPY操作不由外层计划重新分配；候选只覆盖非COPY计算操作，实际输入、输出和跨核COPY由官方构图器根据方案生成。每份候选还须同时通过式\eqref{eq:plan-coverage}和式\eqref{eq:plan-sequence-unique}的覆盖、唯一性检查；这些是候选求解器的格式检查，核级Task、物理版本和事件时间线仍由题给官方评估器展开。对每个逻辑张量$t$，索引保存其大小$b_t$、生产者$p(t)$、去重后的计算消费者$U(t)$、存储区位置以及最终写回指示$o_t$。本研究使用的100张图均按同一字段定义建立这些关系，候选的归属与顺序变化不改变原计算依赖。

问题一按消费Task去重，问题二按消费核心去重。定义计算操作$v$的驻留域$d_B(v)=a(g(v))$，则张量$t$的核级消费域为
\begin{equation}
\mathcal D_t^B=\{d_B(v):v\in U(t)\}.
\label{eq:b-domain}
\end{equation}
同一张量在一个核心上的多个计算消费者，在式\eqref{eq:b-domain}中只占一个元素。这是接收核心去重，不表示张量跨不同核心共享同一片上物理副本。一个核心上的子图即使被分成多个块，其生产与消费边仍保留；官方核内排序根据提交的子图顺序安排这些边所约束的操作。张量的逻辑编号用于统计接收域，后续容量检查则使用换出、恢复形成的物理版本编号，两种索引各有用途。

\begin{table}[H]\centering\caption{问题一与问题二的驻留及同步边界}\label{tab:b-rule-contrast}
\begin{tabularx}{\linewidth}{lYY}\toprule
对象 & 问题一：Task驻留 & 问题二：核心驻留\\\midrule
片上副本可复用范围 & 单个子图Task & 同核合并后的Task\\
相同输入的读入次数 & 按消费子图去重 & 按消费核心去重\\
中间结果的跨域COPY & 源Task写出、各远端Task读入 & 每个远端接收核心各有写出和读入一对\\
跨核释放起点 & 前驱Task完成并等待1000 cycle & 源端COPY\_OUT完成并等待500 cycle\\
容量统计 & Task内部的物理副本 & 同核合并序上的物理副本\\\bottomrule
\end{tabularx}\end{table}

\subsubsection{跨核结构搬运}
问题二允许同核多个子图复用片上数据，因而同一个张量在一个核心上只需一个驻留副本。对输入$t$记$I_B(t)=\{a(g(v)):v\in U(t)\}$；对有生产者的$t$记$R_B(t)=\{a(g(v)):v\in U(t),\,a(g(v))\ne a(g(p(t)))\}$。每个远端接收核需要源核写出与目标核读入，最终输出另行写回，得到
\begin{equation}
Q_B^0=Q_C^0=
\sum_{t\in\mathcal T_{\rm in}}b_t|I_B(t)|
+\sum_{t\in\mathcal T_{\rm prod}}b_t(2|R_B(t)|+o_t).
\label{eq:b-bytes}
\end{equation}
式\eqref{eq:b-bytes}中，图输入没有计算生产者，每个接收核心从DDR读入一次，贡献$b_t|I_B(t)|$。中间张量在生产核心形成，每个远端接收核心均需一份源写出与一份目标读入，贡献$2b_t|R_B(t)|$；原图要求输出的张量另写回一次，贡献$b_t o_t$。等式$Q_B^0=Q_C^0$指同一计划在无Cache与有Cache场景下的结构COPY字节相同；问题三改变读服务来源和完成时刻，不删除这条逻辑搬运。

若一份$b$字节输入由同核两个子图消费，问题一可能读入$2b$，问题二只读入$b$。若一份中间结果的生产者位于核心0、两个消费子图分别位于核心1和2，本问需要两对跨核COPY，即$4b$；将两个消费者放到同一个远端核心后变为$2b$，全部放在生产核心则该项为零。问题一对两个远端Task的同一张量只需源Task写出一次、两个Task各读入一次，即$3b$。这些比较只涉及结构搬运，不包括随后因集中数据而产生的Spill。对于既要跨核发送又要最终写回的张量，官方构图器分别创建跨核COPY与最终输出COPY，式中的$o_t$不能省略。

上述例子不会导出统一的$Q_B^0\le Q_A^0$或相反的不等式，令同一生产张量$t$的外部消费Task数为$r_A=|R_A(t)|$，远端消费核心数为$r_B=|R_B(t)|$，并令$o_t$表示最终写回指示。由式\eqref{eq:a-bytes}和式\eqref{eq:b-bytes}，该张量在两种结构账本中的差值为
\begin{equation}
Q_{B,t}^0-Q_{A,t}^0
=b_t\left(2r_B+o_t-r_A-\mathbf 1\{r_A>0\ \text{或}\ o_t=1\}\right).
\label{eq:ab-structure-difference}
\end{equation}
其中$r_A$按Task计数，包含与生产Task同核但属于不同Task的消费者；$r_B$按远端核心计数且不包含生产核心。右端可为正、负或零；它只比较同一计划下的结构COPY，不含Spill、带宽排队和执行时序。因此，A/B最终方案的搬运量与工期分别按各自账本和官方事件结果评价。

同核切块可以改变执行顺序而不改变式\eqref{eq:b-bytes}的结构字节数；但张量驻留时间和容量恢复仍会变化。随题官方评估程序规定，跨核目标COPY\_IN需等待源COPY\_OUT完成，再加500 cycle释放延迟。

若$e^{\rm out}_{t,i\to j}$和$s^{\rm in}_{t,i\to j}$分别是该边源COPY\_OUT的结束和目标COPY\_IN的开始时刻，则
\begin{equation}
s^{\rm in}_{t,i\to j}\ge e^{\rm out}_{t,i\to j}+500.
\label{eq:b-copy-release}
\end{equation}
源端写出受MTE3流水线及共享DDR服务影响；其完成时间移动，目标读入的可发射时间随之移动。500 cycle属于源端写出完成后的目标读入释放延迟；式\eqref{eq:b-bytes}给出结构字节，式\eqref{eq:objective}给出事件模拟的完成时间。

公式中的去重以消费核心为单位。把同一输入的消费者集中到一核可减少重复读入；把中间张量的消费者移回生产核可减少写读一对COPY。两种操作也可能使一个核心的计算和存活张量过于集中，结构搬运下降与最终工期的关系由完整事件时间线决定。

\subsection{物理副本、容量与Spill}
\subsubsection{物理副本存活期}
片上缓冲的填充、访问和细粒度同步可以作为显式数据编排对象\cite{pellauer2019buffets}。在本题的官方执行规则下，核心$k$接收或生产一份逻辑张量后，该副本从片上空间申请时刻$A_{\tilde t}$开始占用L1或UB，直到最后一个读者完成的时刻$L_{\tilde t}$释放。读者包括计算操作以及将结果写入DDR的COPY。若Step2将它换出并在下次使用前恢复，恢复后占用空间的是新的物理版本$\tilde t'$。对存储区$q$，执行过程必须满足
\begin{equation}
\mathrm{Live}_{kq}(\tau)
=\sum_{\tilde t\in\mathcal V_{kq}(X)}b_{\tilde t}
\mathbf1\{A_{\tilde t}\le\tau<L_{\tilde t}\}
\le C_q,\qquad q\in\{\mathrm{L1},\mathrm{UB}\}.
\label{eq:b-physical-live}
\end{equation}
这里$\mathcal V_{kq}(X)$是计划$X$在核心$k$、存储区$q$产生的物理版本集合，容量$C_q$沿用式\eqref{eq:capacity}。分核与核内顺序改变操作的申请和末次使用次序，因而即使式\eqref{eq:b-bytes}的结构字节相同，式\eqref{eq:b-physical-live}中的重叠区间也会不同。

\begin{figure}[H]\centering
\begin{tikzpicture}[x=0.95cm,y=0.95cm]
\draw[alink] (0,0)--(12.0,0) node[right]{$\tau$};
\draw[very thick,blue!70!black] (1.0,0.9)--(10.3,0.9);
\draw[very thick,orange!80!black] (3.1,1.55)--(8.4,1.55);
\foreach \x in {1.0,3.1,8.4,10.3}
  \draw[black!65] (\x,-0.12)--(\x,0.13);
\node[font=\small,below,align=center] at (1.0,-0.16){第一份\\申请};
\node[font=\small,below,align=center] at (3.1,-0.16){第二份\\申请};
\node[font=\small,below,align=center] at (8.4,-0.16){第二份\\末次读};
\node[font=\small,below,align=center] at (10.3,-0.16){第一份\\末次读};
\node[left,font=\small] at (1.0,0.9){$\tilde t_1$};
\node[left,font=\small] at (3.1,1.55){$\tilde t_2$};
\draw[densely dashed,black!50] (3.1,0.2)--(3.1,1.85);
\draw[densely dashed,black!50] (8.4,0.2)--(8.4,1.85);
\node[font=\small] at (5.8,2.05){两份副本同时占用片上空间};
\end{tikzpicture}
\caption{同核两份物理副本的生存区间及容量重叠}\label{fig:b-lifetime}
\end{figure}

例如，两份各81920 B的UB张量若在图\ref{fig:b-lifetime}所示的区间内同时存活，峰值为163840 B，超过131072 B的UB容量32768 B。前移其中一份的末次使用或推迟另一份的申请可以消除这一重叠；否则Step2须通过换出、恢复等方式组织空间。这一算例只说明容量机制，具体是否发生Spill以官方核内展开为准。

\subsubsection{换出恢复与容量约束}
ZigZag将存储层次与空间、时间映射联合探索\cite{mei2021zigzag}。本题的容量代理只承担候选排序：完整构图之前，程序在每核的局部拓扑操作序上建立位置索引；每个张量在该核上的局部生产者与消费者中，最早操作位置为代理生存起点，最晚操作位置为代理终点。对计划$X$沿这些位置统计张量的理想存活字节$\widehat{\mathrm{Live}}_{kq}(j)$，以
\begin{equation}
\Pi(X)=\sum_{k,q}\left[\max_j\widehat{\mathrm{Live}}_{kq}(j)-C_q\right]_+
\label{eq:pressure}
\end{equation}
描述未换出时的容量压力。式中$[z]_+=\max\{z,0\}$；它对一份张量在某核心只设一段代理区间，不模拟操作级发射、物理版本和内存位置复用。它只在轻量层排列候选，实际容量约束由式\eqref{eq:b-physical-live}及官方展开结果确定。随题官方核内程序依次进行Step1初排及子图优先级调整、Step2插入Spill、Step3补足内存复用依赖；任何一步出现依赖环，候选不进入成功集合。

一次Spill记录的张量大小为$b_s$；恢复COPY\_IN一定搬运$b_s$，若本次换出还需把当前数据写回DDR，则另有一份$b_s$的COPY\_OUT。以官方字段$z_s\in\{0,1\}$表示这次换出是否实际复制数据，Spill字节可从记录逐项复算为
\begin{equation}
Q_B^{\rm spill}=\sum_{s\in\mathcal S_{\rm spill}} b_s(1+z_s).
\label{eq:b-spill}
\end{equation}
该量不是溢出峰值$\Pi$本身：相同容量超额可能因恢复次数、字节数和读者次序不同而产生不同搬运。Spill也可能推迟关键计算的可发射时刻，故工期还取决于它在执行时间线中的位置。

由结构COPY、Spill和原图必要COPY的关系，官方输出中的调度COPY与额外搬运满足
\begin{equation}
Q_B^{\rm sched}=Q_B^0+Q_B^{\rm spill},\qquad
D_B=Q_B^{\rm sched}-Q^{\rm orig}
=(Q_B^0-Q^{\rm orig})+Q_B^{\rm spill}.
\label{eq:b-breakdown}
\end{equation}
第一项是当前分核与核级边界产生的结构附加搬运，第二项是容量恢复搬运。原图、结构及Spill字段均用B计；需要换成MiB时统一除以$2^{20}$。计算工期时还要把流水线占用、跨核释放和共享DDR带宽同时放入事件模拟。下界\eqref{eq:lower-bound}仅解释距离，不参与选优。

为检验式\eqref{eq:b-breakdown}在实际方案中的量级，表\ref{tab:b-copy-ledger}把每个核数下100张图的最终官方结果逐图分解后取算术平均。四列字节量按$2^{20}$ B/MiB换算；未舍入的调度总量等于原图COPY、结构新增与Spill新增三项之和，逐行在B场景的500个结果中均严格成立。

\begin{table}[H]\centering\small
\caption{问题二最终方案的COPY分解（每核数100图；字节列为逐图均值，MiB）}\label{tab:b-copy-ledger}
\begin{tabular}{crrrrr}\toprule
核数 & 原图COPY & 结构新增 & Spill新增 & 调度总量 & 有Spill图数\\\midrule
1 & 2.30 & 0.00 & 24.86 & 24.86 & 81\\
2 & 2.30 & 1.60 & 11.64 & 13.24 & 45\\
3 & 2.30 & 2.50 & 11.39 & 13.89 & 41\\
4 & 2.30 & 3.52 & 10.83 & 14.35 & 38\\
5 & 2.30 & 4.59 & 10.28 & 14.87 & 38\\\bottomrule
\end{tabular}\end{table}

原图COPY均值在五档核数下相同，因为输入图集合没有变化。逐图继承后，由单核到五核，结构新增均值从0升至4.59 MiB，Spill新增均值从24.86降至10.28 MiB；两者共同使调度COPY总量在二核时降至13.24 MiB，五核为14.87 MiB。单核有81张图发生Spill，五核有38张；这些差异来自不同核数下继承的计划和物理驻留过程。搬运字节与工期是两个指标，COPY所处时间、流水线和共享DDR竞争由官方事件模拟决定。

\subsection{初始方案与求解}
\subsubsection{核级驻留方案}
多DNN加速器研究表明，资源分区与层执行调度需要联合考虑\cite{kwon2021heterogeneous}。本题在固定硬件配置下比较分核和核内顺序。整图B0保留原计算图在一个核级Task内的完整驻留，但多核配置中其余核心空闲。B1先求计算依赖的弱连通分量，按分量$\max(M,V)$工作量降序排列，再装入当前桶$\max(M,V)$工作量最小的核心；每个非空核心形成一个子图。无跨分量计算依赖，B1把天然并行的部分分散到不同核心，却无法切开一个占主导的大分量。B2和B3分别使用两种全图就绪序，按剩余$\max(M,V)$工作量目标切成至多$4K$、$8K$个连续块；接近目标的切点优先选择跨切点张量字节较少者。每个新增块接入核心时，对临时计划调用轻量评分，以共同考虑路径、核心负载、结构COPY和容量压力。

\begin{table}[H]\centering\caption{问题二四个起点的实际构造及作用}\label{tab:b-seeds}
\begin{tabularx}{\linewidth}{cYY}\toprule
起点 & 构造规则 & 对核级驻留的影响\\\midrule
B0 & 整图一块，放在核心0 & 无新增跨核边界，保留最长驻留区间\\
B1 & 弱连通分量降序装入最轻核心 & 利用独立分量，子图数量不超过非空核数\\
B2 & 沿H序按$\max(M,V)$目标切至多$4K$块 & 优先考虑长计算链，并在候选切点中压低跨切点字节\\
B3 & 沿R序按$\max(M,V)$目标切至多$8K$块 & 将末次消费释放与新输出规模纳入切点选择\\\bottomrule
\end{tabularx}\end{table}

H序使用式\eqref{eq:a-height}的后继最长计算路径优先级；它在每一步仅从已经就绪的操作中选择。R序在此基础上加入逻辑张量的剩余消费者数。令$n_j(t)$为排入第$j$个操作前张量$t$尚未安排的消费者数，从就绪集合中先取H优先级最高的至多8个，组成$\mathcal R_j^{(8)}$。然后计算
\begin{equation}
r_j(v)=\sum_{\substack{t:\,v\in U(t)\\n_j(t)=1}}b_t
-\sum_{t:\,p(t)=v}b_t,
\qquad v_j^*=\mathop{\arg\max}_{v\in\mathcal R_j^{(8)}}
\bigl(r_j(v),h(v),-\operatorname{id}(v)\bigr).
\label{eq:b-r-order}
\end{equation}
第一项只统计本步成为末次消费者的输入，而非全部输入字节；第二项是本步新产生的输出大小。若多个操作释放量相同，先选H优先级高者，再选编号小者。R序是一种全图静态构造顺序，尚未决定每核物理副本的真正末次读时刻；式\eqref{eq:b-physical-live}的生存区间仍由官方展开确定。

\subsubsection{方案调整与候选筛选}
每个起点保留原方案，并尝试迁块、同核合并、调序和容量风险块的前后移动。问题一同核数的最终方案也作为问题二的热启动锚点，用于直接比较场景变化带来的工期差异。迁块从累计计算周期最大的核心取其列表首块，移至其他核心中最轻者的列表末端；合并保持同核驻留域，却改变提交边界；调序不改变操作归属与式\eqref{eq:b-bytes}，但可改变式\eqref{eq:b-physical-live}中的重叠。每份编辑方案先规范化，再检查操作覆盖、商图无环及同核直接依赖，非法方案跳过；跨起点产生的相同计划按内容去重。

对合法候选，记$L_k(X)$为核心$k$的计算周期总和，$\widehat L_B(X)$为沿原计算拓扑序的路径代理，跨核依赖在代理中附加500 cycle。令$\widehat Q_B(X)$为候选求解器的静态COPY代理，$\Pi(X)$为式\eqref{eq:pressure}，则入围排序采用
\begin{equation}
\widehat F_B(X)=\max\left\{\widehat L_B(X),\max_k L_k(X),
\left\lceil\frac{\widehat Q_B(X)}{60}\right\rceil\right\},\qquad
\widehat\kappa_B(X)=\bigl(\widehat F_B,\widehat Q_B,\Pi\bigr).
\label{eq:b-light-score}
\end{equation}
$\widehat Q_B$按逻辑张量及其接收核心快速累加，计数规则比官方B场景的逐Task展开粗，精确结构字节仍采用式\eqref{eq:b-bytes}。$\widehat L_B$也没有模拟源端COPY耗时、流水线争用或共享DDR，因此三项用于候选排序，最终计划由题给官方评估器返回的$F_B$和$D_B$选择。

\subsubsection{官方评价与方案选定}
官方完整评价先覆盖去重后的整图、分量装桶、上一核已选计划和问题一热启动计划，再从其余候选中至少评价一个代表。设去重后必评计划数为$m$、请求额度为$q$，常规完整调用额度为$\max\{\min(q,2),m+1\}$。题给官方评估器的profile阶段构造核级Task及结构COPY，执行核内Step1/2/3并返回COPY与Spill统计；profile失败时保留错误文本并尝试下一份候选，不消耗完整模拟额度。profile成功后才调用同一评估器的B场景事件模拟，计算每条流水线、DDR及500 cycle跨核释放共同决定的$F_B$。完整模拟即使失败也占一次调用额度；若当前最优慢于上一核数的继承计划，则补评尚未评价的对应回退计划。统一单核基准的回退补评仅在A场景执行。最终在全部成功结果中按$(F_B,D_B,\text{候选名})$取方案；正文逐图B结果再按$1,\ldots,K$核计划与整图基准取最短工期并补齐空核。热启动候选的工期、评价次数和最终去留在受控对照中单独汇总。

\begin{algorithm}[H]\caption{问题二的核级驻留候选与官方评价}\label{alg:b-solve}
\begin{algorithmic}[1]
\REQUIRE 原始图$G$、核心数$K$、官方容量和带宽参数及请求完整调用额度$q$
\ENSURE 最终计划$X_B^*$、工期$F_B$、结构及Spill搬运
\STATE 从原图建立计算依赖、张量生产消费索引、H序与R序
\STATE 生成B0整图、B1分量装桶、B2的$4K$块、B3的$8K$块
\FOR{每个起点}
\STATE 保留原方案；各尝试首块迁核、同核前两块合并与交换
\STATE 检查覆盖和依赖，规范化并按计划内容去重
\ENDFOR
\STATE 按$(\widehat F_B,\widehat Q_B,\Pi,\text{候选名})$排列候选
\STATE 将整图、分量装桶、上一核继承计划和问题一热启动计划置于去重后的必评队列
\STATE 令$m$为必评数，完整调用额度为$\max\{\min(q,2),m+1\}$
\WHILE{尚有候选且未达到常规完整调用额度}
\STATE 对当前候选运行官方profile；若失败，记录错误并继续
\STATE 运行官方B场景事件模拟，保存成功结果或失败文本
\ENDWHILE
\STATE 若当前最优慢于上一核继承计划，补评尚未评价的回退计划
\STATE 在成功集合按$(F_B,D_B,\text{候选名})$选$X_B^*$
\end{algorithmic}\end{algorithm}

伪代码已列出从逻辑张量索引、候选生成、容量筛选到官方事件模拟的步骤。候选求解器使用有限、确定性的起点和局部搜索代表，不对所有边界和操作反复展开。容量压力由操作级区间事件扫描作为候选排序指标，真实换出和恢复由题给官方评估器展开；首个成功候选到最终候选的统计只描述已评价方案之间的选择差异。超过20000个计算操作的五张图也沿用B0--B3构造并纳入100图结果。
\FloatBarrier

\subsection{问题二结果分析}
\subsubsection{逐核工期与逐图结果}
三问使用同一整图单核A0的官方工期$T_{1,i}$。对每张图$i$和目标核数$K$，在已评价的$1,\ldots,K$核方案及整图单核基准中取Makespan最短者，低核方案以空核列表补齐为$K$核方案，再计算$T_{1,i}/F_B(X_{iK}^*)$。因此表\ref{tab:b-results}中各核数均不低于整图单核基准；题面单核曲线仍按定义显示为1。图\ref{fig:speedup-b}的区间来自固定100张图进行1500次有放回实例重采样，展示图集合组成改变时均值的波动，不是芯片重复运行噪声。
\begin{table}[H]\centering\small\caption{问题二：100图逐核工期比与候选改善}\label{tab:b-results}
\begin{tabular}{crrrr}\toprule
核数 & 平均加速比 & 中位数 & 慢于A0的图数 & 初始至最终平均降幅\\\midrule
1 & 1.0000 & 1.0000 & 0 & 1.6528\%\\
2 & 1.7122 & 1.9647 & 0 & 0.3034\%\\
3 & 2.3140 & 2.8340 & 0 & 1.5168\%\\
4 & 2.9183 & 3.6353 & 0 & 2.6911\%\\
5 & 3.3862 & 4.0046 & 0 & 6.3530\%\\\bottomrule
\end{tabular}\end{table}
\begin{figure}[!htbp]\centering
\includegraphics[width=0.91\linewidth,height=0.43\textheight,keepaspectratio]{figures/speedup_B.pdf}
\caption{问题二逐核加速比：均值、中位数、均值95\%实例重采样区间及4至5核重点区间}\label{fig:speedup-b}
\end{figure}

五核平均加速比为3.3862，中位数4.0046，均值的95\%实例重采样区间为$[3.0686,3.6824]$。由四核到五核，均值从2.9183升至3.3862，增加0.4679；同图配对中65张图工期缩短、35张图保持不变，没有图因增加核数而变慢。表中“初始至最终平均降幅”逐图计算$(F_{\rm init}-F_{\rm final})/F_{\rm init}$后取均值；五核为6.3530\%，其中21张图的最终候选缩短工期。该指标衡量已评候选的选择收益。

\subsubsection{A、B场景结构搬运对比}
逐图继承后的五核B方案平均额外搬运为15596420.30 B。按式\eqref{eq:b-breakdown}分解，结构附加部分为4814112.38 B，Spill为10782307.92 B；100张图中38张出现非零Spill。平均值受少数大图和多次恢复拉动，容量恢复只主导其中部分图；最终工期由继承后的完整事件时间线确定。

两场景在同一100张图、五核终态上的平均加速比分别为A的3.4132与B的3.3862；逐图比较为B快、相同或变慢的数量取决于两套场景的计划和执行规则。本文将逐图继承后的A/B结果用于总体曲线，将固定五核计划的同图对照保留在受控消融中；两种口径分别回答“选核后的最终工期”和“同一核数计划的场景差异”。图\ref{fig:b-ab-tradeoff}沿用固定五核计划配对，将每张图的额外搬运差放在横轴、B/A工期比放在纵轴，并标出方向相反的代表图。

问题二相对问题一的均值差同时包含执行规则和所选计划的变化。核级复用减少同一核心内多个消费者的重复读入；核级 Task 也可能延长物理副本的存活区间，增加 L1、UB 中的区间重叠，跨核中间结果仍需源端写出和目标端读入。由此节省的结构搬运可能被容量恢复、跨核同步和关键路径位置抵消。A、B的五核均值仅相差0.0091，逐图却同时出现改善、持平与退化。固定计划对照用于分离执行规则效应，结构搬运、Spill 和工期则共同描述各图最终方案。
\samplefigure{ab_tradeoff_scatter}{五核逐图A/B对照：额外搬运差与工期比；两轴按图内刻度解释}{fig:b-ab-tradeoff}

\begin{table}[H]\centering\small\caption{五核代表图在A/B场景的工期及额外搬运分解}\label{tab:b-cases}
\begin{tabular}{llrrr}\toprule
图 & 场景及候选 & 工期/cycle & 结构附加/B & Spill/B\\\midrule
\Case{055} & A：A2\_move & 110002 & 199780 & 0\\
 & B：B1\_base & 25635 & 231424 & 0\\
\Case{053} & A：A1\_base & 492699 & 3104256 & 3814656\\
 & B：B0\_base & 1472686 & 0 & 5405536\\\bottomrule
\end{tabular}\end{table}

图\ref{fig:b-case055-timeline}进一步展开\Case{055}五核B最终方案的官方事件时间线。上半部以浅色条表示五个核级Task的起止区间，而非计算流水线始终忙碌；蓝线、橙线分别画出各核COPY\_IN与COPY\_OUT的历时，橙色菱形标出COPY\_OUT的起点。下半部由同一批COPY事件的开始和结束时刻，精确累计同时在途的DDR COPY数。该图使用官方结果中的252次COPY事件，不是按平均字节量合成的示意轨迹。
\samplefigure{b_case055_official_timeline}{\Case{055}五核B最终方案的核级Task、COPY及在途搬运时间线}{fig:b-case055-timeline}

五个Task均从零时刻开始，最后完成的是核心1的Task，时刻为25635 cycle；前段同时在途COPY最多为6次，约14000 cycle之后搬运明显稀疏。图示说明共享DDR的并发搬运数随时间变化，总搬运量与事件时间线共同决定搬运延迟。它只展示B最终方案内部的执行结构，A/B工期差还需结合两种场景各自的完整事件结果判断。

\Case{055}的B工期仅为A的约0.233倍，结构附加搬运却从199780 B增加到231424 B，两者均无Spill。更快方案并非简单减少了字节，而是改变了计划与Task执行规则后缩短了决定完成时间的部分。\Case{053}呈相反方向：B选中整图B0，结构附加为零，但Spill达到5405536 B，工期比A增加到约2.99倍；A虽有3104256 B结构附加和3814656 B Spill，仍更早完成。B0在该图上集中计算，零结构通信不足以抵消失去多核并行及较长驻留区间的综合代价。这两例说明式\eqref{eq:b-bytes}与式\eqref{eq:b-physical-live}必须共同进入评价，而真实胜负由事件模拟决定。

\subsubsection{容量恢复算例}
对于同一场景的两份已评候选，保持结构附加搬运相同可以单独观察块结构与核内次序的综合影响。\Case{008}单核B3\_base和B3\_aggregate的结构附加均为0，前者Spill为10562688 B、工期500608 cycle，后者Spill为7059072 B、工期405761 cycle。合并改变块结构及官方优先次序，Spill减少3503616 B，工期缩短94847 cycle。这个候选显示相同逻辑结构字节数仍会出现不同容量恢复与完成时间；差异同时来自块合并和排序变化。

\begin{table}[H]\centering\caption{\Case{008}单核问题二两份成功候选}\label{tab:b-case008}
\begin{tabular}{lrrr}\toprule
候选 & 工期/cycle & 结构附加/B & Spill/B\\\midrule
B3\_base & 500608 & 0 & 10562688\\
B3\_aggregate & 405761 & 0 & 7059072\\\bottomrule
\end{tabular}\end{table}

从候选求解器与题给评估器的记录看，B主矩阵对每个图与核数均保存首次成功和最终计划、候选profile、正式结果及失败文本。复核时先验证操作覆盖和商图/核内直接依赖，再核对官方报告中的跨核COPY释放、L1与UB峰值以及式\eqref{eq:b-breakdown}的字节恒等式。容量峰值核对的是官方展开报告与配置的一致性；物理版本在整个事件时间线的逐时刻重建另需独立轨迹检查。上述结果对应有限候选的官方事件模拟评价；运行平台为CPU官方Python评估器。

\subsection{问题二结论}
问题二的核级接收域把同核消费归并到一份结构读入，式\eqref{eq:b-bytes}给出跨核COPY账本，式\eqref{eq:b-physical-live}及式\eqref{eq:b-spill}解释容量恢复为何会改变工期。按逐图继承规则，五核平均加速比为3.3862，所有图均不慢于整图单核基准；每张图的选中结果保留计划和事件时间线。

提交计划包含操作映射\texttt{node\_to\_subgraph}和核心列表\texttt{core\_schedules}。逐图保存的最终计划是问题三的固定B起点$X_B^*$；选择记录与候选结果分别保存工期、搬运和计划身份。问题三在相同计划上加入只读Cache，再比较缓存硬件作用与方案适配。B计划的哈希在500组配对中逐组核对。

\begin{table}[H]\centering\small\caption{问题二可追溯的方案与评价记录}\label{tab:b-artifacts}
\begin{tabularx}{\linewidth}{lYY}\toprule
记录 & 内容 & 在下一问中的用途\\\midrule
\texttt{final\_plan.json} & 操作映射、分核和同核列表 & 固定B计划及C的初始方案\\
\texttt{final\_selection.json} & 候选名、计划哈希、工期、搬运 & 核对B计划身份\\
候选\texttt{profile.json} & 核级展开、COPY及Spill统计 & 解释片上恢复和候选入围\\
候选\texttt{result.json} & 工期、流水线事件与搬运字段 & 同计划B/C配对及约束核验\\\bottomrule
\end{tabularx}\end{table}

\subsection{问题二灵敏度分析}
问题二的片上驻留对分块方式敏感。保持\Case{008}、单核和B场景硬件配置不变，比较同一拓扑切块起点与相邻块合并后的两份已评候选：表\ref{tab:b-case008}中结构附加搬运均为0，Spill由10562688 B降至7059072 B，工期由500608降至405761 cycle。分块调整使恢复搬运减少3503616 B，同时使工期缩短94847 cycle；这两个变化来自同一组官方结果，不把Spill字节降幅直接等同于工期降幅。

该单图扰动说明，核级Task内的物理副本生存期会随块边界和操作次序改变。表\ref{tab:b-results}和图\ref{fig:speedup-b}给出100图的逐核结果：五核均值为3.3862，四至五核配对中65图工期缩短、35图不变。固定硬件下的块结构调整与逐图继承增加核心数分别改变不同的执行条件，须结合各自的工期和搬运结果分析。
